\par\bigskip
\begin{center}
\rule{.88\textwidth}{.3pt}\par
Göttingische gelehrte Anzeigen.\quad 25. September 1826.\par
\rule{.88\textwidth}{.3pt}
\end{center}
\pdfbookmark[1]{25. September 1826}{anzeige-1826-09-25}

Am 16. September überreichte Herr Hofrat \textsc{Gauss} der Königlichen Societät eine Vorlesung mit dem Titel:
\begin{center}
\textit{Supplementum Theoriae combinationis observationum erroribus minimis obnoxiae.}
\end{center}
\noindent Allen früheren Arbeiten zur Anwendung der Wahrscheinlichkeitsrechnung auf die zweckmäßigste Nutzung von Beobachtungen, besonders auch der Behandlung dieses Gegenstands im fünften Band der \textit{Commentationes recentiores}, liegt hinsichtlich der Form der Hauptaufgabe eine bestimmte Voraussetzung zugrunde, die den meisten in der Praxis vorkommenden Fällen durchaus entspricht. Sie besteht darin, dass die be\SourceLeaf{110}%
obachteten Größen in bekannter Weise von gewissen unbekannten Größen (Elementen) abhängen, das heißt bekannte Funktionen dieser Elemente sind. Damit die Aufgabe überhaupt zu dieser Klasse gehört, muss die Anzahl der Elemente kleiner sein als die der beobachteten Größen, sodass die beobachteten Größen ihrerseits voneinander abhängig sind.

Dennoch ist die angenommene Voraussetzung in vielen Fällen nicht unmittelbar erfüllt: Die beobachteten Größen sind noch nicht in der Form bekannter Funktionen gewisser unbekannter Elemente gegeben, und man sieht nicht sofort, wie sie sich in diese Form bringen lassen. Stattdessen wird ihr gegenseitiger Zusammenhang – der natürlich irgendwie gegeben sein muss – durch bestimmte Bedingungsgleichungen ausgedrückt, denen die wahren Werte dieser Größen der Natur der Sache nach notwendig genau genügen müssen. Bei näherer Betrachtung erkennt man bald, dass sich dieser Fall nicht wesentlich vom anderen unterscheidet, sondern nur in der Form, und sich theoretisch tatsächlich leicht auf jenen zurückführen lässt. In der Anwendung bleibt dies jedoch häufig ein unnatürlicher Umweg, der erheblich beschwerlichere Rechnungen mit sich bringt als eine eigene, der ursprünglichen Gestalt der Aufgabe besonders angepasste Lösung. Eine solche Lösung ist Gegenstand der vorliegenden Abhandlung. Die Lösung der Aufgabe, die die Abhandlung als selbstständiges, von der früheren Abhandlung unabhängiges Ganzes gibt, ist ihrerseits so geschmeidig, dass es in manchen Fällen sogar vorteilhaft sein kann, sie dort anzuwenden, wo die der älteren Methode zugrunde liegende Voraussetzung bereits erfüllt ist.

Die Hauptaufgabe stellt sich nun folgendermaßen dar. Wenn eine andere Größe in irgendeiner Weise von den Größen \(v,\,v',\,v''\) usw. abhängt, zwischen denen ein durch eine oder mehrere Bedingungsgleichungen gegebener Zusammenhang besteht, und zum Beispiel durch die Funktion \(u\) ausgedrückt werden kann, lässt sie sich ebenso auf unendlich viele andere Arten aus jenen Größen bestimmen oder statt durch \(u\) durch unendlich viele andere Funktionen ausdrücken. Selbstverständlich ergeben alle diese Funktionen dasselbe Resultat, sofern man die wahren Werte von \(v,\,v',\,v''\) usw. einsetzt, die allen Bedingungsgleichungen genügen. Hat man jedoch nur Näherungswerte von \(v,\,v',\,v''\) usw., wie sie Beobachtungen mit beschränkter Genauigkeit stets nur liefern können, so \SourceLeaf{111}%
können auch die daraus abgeleiteten Größen keine absolute Richtigkeit beanspruchen. Im Allgemeinen werden die verschiedenen für \(u\) verwendeten Funktionen ungleiche, vor allem aber ungleich zuverlässige Ergebnisse liefern. Die Aufgabe besteht darin, aus der unendlichen Mannigfaltigkeit der Funktionen, durch die sich die unbekannte Größe ausdrücken lässt, diejenige auszuwählen, bei deren Ergebnis die kleinstmögliche Unzuverlässigkeit zu befürchten ist.

Die Abhandlung gibt eigentlich \textit{zwei} Lösungen dieser Aufgabe. Die erste erreicht das Ziel auf dem kürzesten Weg, wenn tatsächlich nur eine unbekannte, in vorgeschriebener Weise von den Beobachtungen abhängige Größe abzuleiten ist. Die nähere Betrachtung dieser Lösung führt zugleich zu dem bemerkenswerten Satz, dass man für die unbekannte Größe genau denselben Wert erhält wie aus der zweckmäßigsten Kombination der Beobachtungen, wenn man an den Beobachtungen gewisse nach bestimmten Regeln berechnete Änderungen vornimmt und sie dann in eine beliebige Funktion einsetzt, welche die unbekannte Größe ausdrückt. Diese Änderungen erfüllen nicht nur alle Bedingungsgleichungen; unter allen denkbaren Systemen, die dies tun, haben sie auch die Eigenschaft, dass die Summe ihrer Quadrate – sofern die Beobachtungen als gleich zuverlässig vorausgesetzt werden – möglichst klein ist. Damit ergibt sich zugleich eine neue Begründung der Methode der kleinsten Quadrate. Diese von der Funktion \(u\) völlig unabhängige \textit{Ausgleichung} der Beobachtungen stellt eine zweite Lösungsart dar, die gegenüber der ersten einen großen Vorzug hat, wenn mehr als eine unbekannte Größe auf zweckmäßigste Weise aus den Beobachtungen abzuleiten ist: Dadurch werden die Beobachtungen im Voraus für \textit{jede} von ihnen zu machende Anwendung vorbereitet. Bei dieser zweiten Lösung musste noch eine besondere Anleitung hinzukommen, wie sich der bei jeder einzelnen Anwendung erreichte Genauigkeitsgrad bestimmen lässt. Zu all dem enthält die Abhandlung vollständige und möglichst einfache Vorschriften, die sich natürlich hier nicht auszugsweise wiedergeben lassen. Ebenso können wir uns hier nicht auf Einzelheiten der weiteren Untersuchungen einlassen, die nach der Entwicklung der Hauptaufgaben ausgeführt wurden und eng mit dem Gegenstand zusammenhängen. Nur einen bemerkenswerten Satz führen wir an: \SourceLeaf{112}%
Die Vorschriften zur vollständigen Ausgleichung der Beobachtungen ergeben stets dasselbe Resultat, ob sie auf die ursprünglichen Beobachtungen selbst oder auf Beobachtungen angewandt werden, die vorläufig bereits \textit{unvollständig} ausgeglichen wurden – sofern man diesen Begriff in der näher bestimmten Bedeutung der Abhandlung versteht. Als Spezialfall ist darunter auch derjenige begriffen, in dem mit den Beobachtungen bereits eine zwar vorschriftsmäßig durchgeführte, aber nur einen Teil der Bedingungsgleichungen berücksichtigende Ausgleichung vorgenommen worden war.

Der letzte Teil der Abhandlung besteht aus einigen sorgfältig ausgearbeiteten Beispielen für die Anwendung der Methode. Sie stammen teils aus den geodätischen Messungen des Generals von \textsc{Krayenhoff}, teils aus der vom Verfasser selbst im Königreich Hannover durchgeführten Triangulierung. Sie können sowohl die Anwendung dieser Theorie weiter erläutern als auch manche Umstände solcher Messungen überhaupt in helleres Licht rücken.

Die trigonometrischen Messungen gehören ganz besonders in das Gebiet, in dem die Wahrscheinlichkeitsrechnung Anwendung findet, und insbesondere in derjenigen Form, die in der vorliegenden Abhandlung entwickelt wird. Gerade hier ist es die Regel, mehr zu messen, als unbedingt notwendig ist, sodass die Messungen einander vielfach kontrollieren. Nur wenn man die strengen Grundsätze der Wahrscheinlichkeitsrechnung anwendet, kann man aus diesem Umstand den möglichen Vorteil voll ziehen und den Ergebnissen die größte Genauigkeit geben, deren sie fähig sind. Außerdem bieten diese Grundsätze ein Mittel, sowohl die Genauigkeit der Messungen selbst als auch die Zulässigkeit der darauf gegründeten Resultate zu bestimmen. Schließlich dienen sie dazu, unter mehreren möglichen Anordnungen eines Dreieckssystems die zweckmäßigste auszuwählen. All dies geschieht nach festen, zuverlässigen Regeln unter Ausschluss jeder Willkür. Doch sind eine sichere Beurteilung und die bestmögliche Nutzung der Messungen nur möglich, wenn sie in unverfälschter Authentizität und Vollständigkeit vorliegen. Daher wäre sehr zu wünschen, dass alle umfangreicheren Messungen dieser Art, die besondere Genauigkeit beanspruchen, stets mit der erforderlichen Ausführlichkeit veröffentlicht würden. Allzu häufig geschieht das Gegenteil, und es werden nur die Endresultate der einzelnen gemesse\SourceLeaf{113}%
nen Winkel mitgeteilt. Sind solche Endresultate nach richtigen Grundsätzen gebildet worden, indem sämtliche einzelnen Beobachtungsreihen beitragen durften, die nicht mit Sicherheit einen völlig unzulässigen Fehler enthalten, ist der Nachteil zwar bei weitem nicht so groß wie dann, wenn man nur diejenigen Reihen beibehält, die am besten zu den naheliegenden Prüfmitteln passen: den Winkelsummen in jedem Dreieck und den Summen der Horizontalwinkel um jeden Punkt. Wird dieses völlig verwerfliche Verfahren angewandt – sei es aus Unkenntnis der wahren Grundsätze einer richtigen Theorie oder aus dem geheimen Wunsch, den Messungen den Anschein größerer Genauigkeit zu verleihen –, geht der Maßstab für eine gerechte Beurteilung der Beobachtungen und der daraus abzuleitenden Ergebnisse verloren. Die übliche Prüfung anhand der Winkelsummen in einzelnen Dreiecken und an den Punkten, an denen die gemessenen Winkel den ganzen Horizont umfassen, scheint dann eine Messgenauigkeit zu beweisen, die möglicherweise weit davon entfernt ist, tatsächlich erreicht zu werden. Gibt es weitere Prüfmittel durch Seitenverhältnisse in geschlossenen Polygonen oder durch Diagonalrichtungen, so verraten diese mit Gewissheit, dass viel größere Fehler vorhanden sind. Umgekehrt muss der Versuch, die Beobachtungen hinsichtlich der Prüfmittel ohne die zuverlässigen Vorschriften der Wahrscheinlichkeitsrechnung auszugleichen, wenn die zuletzt genannte Voraussetzung gilt, immer ein Herumtappen im Dunkeln bleiben; er führt zu größeren, oft viel größeren Korrekturen, als nötig sind. Dadurch kann leicht ein \textit{zu} ungünstiges Urteil über die Messungen entstehen. Diese Bemerkungen zeigen, wie wichtig sowohl eine hinreichend ausführliche Veröffentlichung als auch eine auf strengen Grundsätzen beruhende mathematische Kombination geodätischer Messungen ist. Sie gelten aber offenbar mehr oder weniger für Beobachtungen jeder Art – astronomische, physikalische und so weiter –, die sich auf quantitative Größen beziehen, sofern die Vielfalt der dabei auftretenden Umstände Mittel zu wechselseitigen Kontrollen bietet.
\par\medskip\centerline{\rule{.3\linewidth}{.3pt}}
